\section{Why Not Accuracy}
\label{sec:eval-why}

Equation~\ref{eq:naive} established that a constant classifier scores 66.95\% on this
dataset. Any metric that a useless model can score well on is not a metric that
distinguishes useful models from each other. Accuracy is reported in the results tables
because omitting it would look evasive, and it is used to decide nothing.

The metrics that do the work are per-class, and two classes carry the weight. Melanoma
is the class whose misdiagnosis is lethal. Basal cell carcinoma is the most common
malignancy in the set. Together they are 16.2\% of the data and essentially all of the
clinical risk.

\section{Metrics}
\label{sec:eval-metrics}

For each class $c$, with $TP_c$, $FP_c$ and $FN_c$ the true positives, false positives
and false negatives:

\begin{align}
  \mathrm{precision}_c &= \frac{TP_c}{TP_c + FP_c} \label{eq:precision}\\[0.4em]
  \mathrm{recall}_c    &= \frac{TP_c}{TP_c + FN_c} \label{eq:recall}\\[0.4em]
  \mathrm{F1}_c        &= 2 \cdot \frac{\mathrm{precision}_c \cdot \mathrm{recall}_c}
                                       {\mathrm{precision}_c + \mathrm{recall}_c}
                          \label{eq:f1}
\end{align}

The asymmetry between the two error types is the whole reason precision and recall are
reported separately rather than only as their harmonic mean. For melanoma, $FN$ is a
cancer sent home as benign. $FP$ is a healthy patient referred for a second look. In a
screening context where a dermatologist reviews the flagged cases, these are not
comparable costs, and a single F1 figure conceals which way a given configuration is
erring.

Three aggregate metrics are also reported. Macro F1 is the unweighted mean of per-class
F1, giving each class equal weight regardless of frequency, which is what makes it
informative here where accuracy is not. Balanced accuracy is the mean of per-class
recall. Macro AUC-OvR is the one-versus-rest area under the ROC curve averaged over
classes, computed from predicted probabilities rather than hard labels, and is therefore
insensitive to the decision threshold.

A $7 \times 7$ confusion matrix is reported per run. Aggregates say how much a model got
wrong; the confusion matrix says what it confused with what, which is the part a
clinician would ask about.

\section{Selection Criteria}
\label{sec:eval-criteria}

The criteria fixed in Chapter~3 are restated here in the form the analysis code
implements.

For RQ2, the score for each run is the unweighted mean of recall on the two clinical
classes, and the best configuration per architecture is the arg-max:

\begin{lstlisting}[caption={The RQ2 criterion, \texttt{scripts/build\_tables.py}.},
                   label={lst:rq2}]
def rq2_score(exp: str) -> float:
    return (report(exp, "mel", "recall") + report(exp, "bcc", "recall")) / 2
\end{lstlisting}

For RQ1, the winners under that criterion are compared on melanoma F1 against the 0.05
threshold. Because ``best'' is a definitional choice, Chapter~7 also reports the
comparison under the alternative reading where best means highest macro F1. If the two
readings disagree about whether RQ1 passes, that is a finding about the fragility of the
criterion and is reported as one.

RQ1 is computed over the confirmatory experiments only, which the code enforces rather
than leaves to discipline:

\begin{lstlisting}[caption={RQ1 restricted to the confirmatory factorial,
\texttt{scripts/build\_tables.py}.}, label={lst:rq1-scope}]
# RQ1 is defined over the confirmatory pair only. B3 is an ablation and does
# not replace either incumbent in the pre-registered question.
conf = [e for e in exps if e in CONFIRMATORY_EXPERIMENTS]
mob = [e for e in conf if EXPERIMENTS[e]["architecture"] == "mobilenet_v2"]
res = [e for e in conf if EXPERIMENTS[e]["architecture"] == "resnet50"]
\end{lstlisting}

\section{Frequency Tiers}
\label{sec:eval-tiers}

Per-class metrics across seven classes are detailed but hard to read as a whole. As a
secondary view, classes are grouped into a majority tier (nv, bkl, mel) and a minority
tier (akiec, vasc, df) by training frequency, and mean F1 is reported for each.

This is a compression of information already present in the per-class table, included
because it makes one specific question legible: whether a configuration is buying
minority-class performance at the expense of the majority. A run whose minority-tier
mean exceeds its majority-tier mean has inverted the usual ordering, which is a strong
signal about what the imbalance remedy did.

\section{Efficiency Measurement}
\label{sec:eval-efficiency}

RQ3 quantities are measured, not quoted. Parameter count is the sum over parameters with
\texttt{requires\_grad}. Model size is the serialised \texttt{state\_dict} in megabytes.
Latency is wall-clock time for a single forward pass at batch size 1, measured on CPU
after ten warm-up iterations, reported as mean and standard deviation over 100 timed
runs.

CPU rather than GPU because the deployment scenario motivating this thesis is an offline
device with no accelerator. A GPU latency figure would be easy to obtain and would not
speak to the question.

EfficientNet-B3 is measured at both 224 and 300 pixels, since latency scales with input
area and a single figure would misrepresent it.

\begin{lstlisting}[caption={Latency measurement, \texttt{scripts/benchmark\_efficiency.py}.},
                   label={lst:latency}]
dummy = torch.randn(1, 3, image_size, image_size, device=device)
times_ms = []
with torch.no_grad():
    for _ in range(WARMUP_RUNS):
        model(dummy)
    for _ in range(TIMED_RUNS):
        start = time.perf_counter()
        model(dummy)
        times_ms.append((time.perf_counter() - start) * 1000)
\end{lstlisting}

Timing on shared or virtualised hardware is noisy, and the measurement host is recorded
alongside the numbers so the figures can be read with that in mind. Where the standard
deviation is large relative to the mean, the comparison between architectures remains
informative because all are measured under identical conditions, but the absolute values
should not be treated as a specification.

\section{Statistical Testing}
\label{sec:eval-stats}

\subsection{McNemar's test}

Following Dietterich~\cite{dietterich1998tests}, comparisons use McNemar's paired test.
For two models evaluated on the same test set, the $2 \times 2$ contingency table counts
images by whether each model classified them correctly:

\begin{table}[htbp]
\centering
\caption{McNemar contingency table. Only $b$ and $c$ enter the test.}
\label{tab:mcnemar-table}
\small
\begin{tabular}{@{}lcc@{}}
\toprule
& \textbf{B correct} & \textbf{B wrong} \\
\midrule
\textbf{A correct} & $a$ & $b$ \\
\textbf{A wrong}   & $c$ & $d$ \\
\bottomrule
\end{tabular}
\end{table}

The null hypothesis is $b = c$: the two models disagree equally often in both
directions. Cells $a$ and $d$ are discarded, which is the essential property. Images both
models handle identically carry no information about which is better, and on a dataset
where two thirds of images belong to an easy majority class, a test that counted them
would be dominated by agreement.

The exact binomial form is used when $b + c < 25$, where the chi-squared approximation is
unreliable, and the continuity-corrected chi-squared form otherwise. The selection is
automatic rather than chosen per comparison:

\begin{lstlisting}[caption={Test variant selection, \texttt{scripts/statistical\_tests.py}.},
                   label={lst:mcnemar}]
discordant = a_only + b_only
# Exact binomial when the discordant count is small; the chi-squared
# approximation is unreliable there.
use_exact = discordant < 25
result = mcnemar([[both_right, a_only], [b_only, both_wrong]],
                 exact=use_exact, correction=not use_exact)
\end{lstlisting}

Predictions are joined on \texttt{image\_id} rather than by row position, and the test
asserts that both runs agree on ground truth before proceeding. A mismatch there would
indicate the two runs saw different splits, which would invalidate the pairing.

\subsection{Two families}
\label{sec:stat-families}

Twenty-one comparisons are declared, in two families corrected independently.

The \textbf{confirmatory family} contains nine comparisons over E1 to E6: three
cross-architecture within strategy, and six cross-strategy within architecture. These
were specified before any results existed.

The \textbf{exploratory family} contains twelve comparisons involving EfficientNet-B3:
six against the incumbents at matched resolution, three across strategies within B3 at
224, and three resolution pairs.

\begin{lstlisting}[caption={The two declared families, \texttt{scripts/statistical\_tests.py}.},
                   label={lst:families}]
CONFIRMATORY = [
    # architecture held constant within strategy
    ("E1", "E2"), ("E3", "E4"), ("E5", "E6"),
    # strategy varied within architecture
    ("E1", "E3"), ("E1", "E5"), ("E3", "E5"),
    ("E2", "E4"), ("E2", "E6"), ("E4", "E6"),
]

EXPLORATORY = [
    # B3 at 224 against each incumbent, strategy held constant
    ("E7", "E1"), ("E7", "E2"),
    ("E8", "E3"), ("E8", "E4"),
    ("E9", "E5"), ("E9", "E6"),
    # strategy varied within B3 at 224
    ("E7", "E8"), ("E7", "E9"), ("E8", "E9"),
    # resolution: same architecture, same strategy, 224 against 300
    ("E7", "E10"), ("E8", "E11"), ("E9", "E12"),
]
\end{lstlisting}

Holm-Bonferroni correction is applied within each family at $\alpha = 0.05$. Holm is
used rather than plain Bonferroni because it is uniformly more powerful while making no
additional assumptions.

Correcting separately is a decision with a consequence, and it should be visible rather
than buried. Pooling all twenty-one into one family would apply a harsher correction to
the nine pre-registered comparisons, so a decision taken after the confirmatory runs had
completed could determine whether those earlier findings reach significance. Separate
families prevent that. The cost is that the exploratory family's own error rate is
controlled only within itself, which is the appropriate treatment for comparisons that
were not planned in advance.

The implementation refuses to correct over a partial family:

\begin{lstlisting}[caption={Refusing partial correction, \texttt{scripts/statistical\_tests.py}.},
                   label={lst:partial}]
    """Holm-correct within one family.

    Skips the family entirely if any member is missing rather than silently
    correcting over a subset: dropping comparisons changes the correction and
    would quietly make the surviving ones look more significant than they are.
    """
\end{lstlisting}

\subsection{What a non-significant result means here}

A comparison that fails to reach significance is not a null result to be passed over. It
states that the two configurations disagree about equally often in both directions. When
one of them is roughly a tenth the size of the other, that is a substantive finding about
deployability rather than an absence of one, and Chapter~7 reports non-significant
comparisons on the same footing as significant ones.
